\section{XSS-Analyse und Behebung}
\label{sec:xss}

\subsection{Gefundene Stored-DOM-XSS-Schwachstelle}

Im eigenen Frontend-Code wurden keine direkten produktiven
Verwendungen von \texttt{v-html}, \texttt{innerHTML},
\texttt{insertAdjacentHTML} oder \texttt{eval} gefunden. Die
relevante Senke lag in der Drittbibliothek \texttt{vue-cal} in
Version \texttt{\^{}4.8.1}. Deren Standardrenderer behandelt
\texttt{event.title} und \texttt{event.content} als HTML.

Mehrere Scheduler-Ansichten übergaben datenbankbasierte Werte an
genau diese Felder, ohne einen eigenen Event-Slot zu definieren.
Besonders relevant waren folgende Komponenten:

\begin{itemize}
  \item \texttt{ReserveTimeslot.vue}: Reservierungslabel und
    Veranstaltungsname als \texttt{title},
  \item \texttt{AssignmentView.vue}: Veranstaltungs- und Gruppenname
    als \texttt{title} beziehungsweise \texttt{content},
  \item \texttt{EditIncompleteStateAvailabilities.vue}: Gruppenname
    als \texttt{content},
  \item \texttt{TimeSlotCalendarView.vue}: Verwendung des
    Standardrenderers und HTML-Strings für Statusicons.
\end{itemize}

Die bereits vorhandene Datei \texttt{ScheduleOverview.vue} zeigte das
sichere Gegenmuster: Sie definierte eigene \texttt{event}-Slots und
gab Titel, Inhalte und Annotationen über Vue-Interpolation aus.

\subsection{Konkreter Datenfluss und Risiko}

Der nachgewiesene Datenfluss verlief über eine persistente Quelle und
eine DOM-Senke:

\begin{center}
  \texttt{Eingabefeld / Import} $\rightarrow$ \texttt{Datenbank}
  $\rightarrow$ \texttt{API} $\rightarrow$
  \texttt{event.title/content} $\rightarrow$ \texttt{vue-cal innerHTML}
\end{center}

Ein harmloser lokaler Nachweis verwendete einen Wert der Form:

\begin{lstlisting}[caption={Lokaler XSS-Demo-Payload}]
XSS Demo <img src=x onerror="alert('XSS Demo')">
\end{lstlisting}

Das nicht vorhandene Bild löst den \texttt{onerror}-Handler aus. Der
Alert beweist lediglich Codeausführung. Realistisch könnte
eingeschleuster Code sichtbare Planungsdaten lesen, UI-Inhalte
manipulieren oder API-Requests im Kontext des eingeloggten Nutzers
auslösen. Das Risiko wird dadurch erhöht, dass
\texttt{web/src/App.vue} die Keycloak-Instanz als
\texttt{window.\$keycloak} verfügbar macht und der Access Token
zusätzlich im Vuex-Session-State liegt.

\subsection{Struktureller Fix}

Die Schwachstelle wurde nicht durch eine Payload-Blacklist behoben,
sondern durch Entfernung der unsicheren Ausgabesenke. Alle relevanten
\texttt{VueCal}-Instanzen mit nutzer- oder datenbankkontrollierten
Werten besitzen zum Berichtsstand eigene Event-Slots. Titel, Inhalte
und Statusinformationen werden darin als Text oder als strukturierte
UI-Elemente gerendert.

Der Demo-Stand \path{origin/its-xss-cors-security-break}
demonstriert den Unterschied exemplarisch mit zwei Komponenten:

\begin{itemize}
  \item \texttt{ReserveTimeslot.vue}: ursprünglicher Standardrenderer
    für die unsichere Demo,
  \item \texttt{ReserveTimeslotSafe.vue}: eigener
    \texttt{v-slot:event} mit \texttt{\{\{ event.title \}\}}.
\end{itemize}

\path{web/.env.xss-safe}, die Umschaltung in \path{web/src/router.js}
und die npm-Skripte \path{dev:xss-safe} beziehungsweise
\path{build:xss-safe} sind reine Demo-Artefakte. Auf
\texttt{capstone-main} ist der sichere Renderer direkt in den
regulären Ansichten umgesetzt; eine Laufzeitumschaltung zurück zur
anfälligen Variante ist dort nicht vorhanden. Die Umsetzung ist in
Commit \texttt{3742df35} mit der Nachricht
\texttt{fix: render VueCal events safely} dokumentiert.

\subsection{Bewertung nach dem Rollout}

Nach dem vollständigen Rollout und einer erneuten Prüfung aller
relevanten Kalenderansichten auf \texttt{capstone-main} besteht dort
kein bekanntes aktuelles XSS-Risiko über den
\texttt{vue-cal}-Standardrenderer. Sämtliche anderen zum
Berichtsstand geprüften Remote-Branches enthalten Commit
\texttt{3742df35} nicht und weisen das ursprüngliche anfällige
VueCal-Muster weiterhin auf. Diese Aussage bedeutet nicht, dass
zukünftige Änderungen automatisch sicher bleiben. Neue Bibliotheken,
neue Kalenderansichten oder eine spätere Verwendung von
\texttt{v-html} können dieselbe Schwachstellenklasse erneut
einführen.

\subsection{Verbleibende Schutzmaßnahmen}

Für nachhaltigen Schutz sind folgende Ergänzungen sinnvoll:

\begin{itemize}
  \item Security-Regressionstests, die typische Payloads in jeder
    Kalenderansicht als Text erwarten,
  \item serverseitige Längen- und Fachvalidierung für Namen, Labels,
    Annotationen und Freitextfelder,
  \item eine produktionsgerechte CSP ohne unnötiges \texttt{unsafe-inline},
  \item regelmäßige Prüfung von Dependency-Updates und Bibliotheksrenderern,
  \item Entfernung unnötiger globaler Token-Zugriffe wie
    \texttt{window.\$keycloak}.
\end{itemize}
